% =====================================================================
\chapter{Contrastes de hipótesis}
\label{tema:5}
% =====================================================================

Un contraste de hipótesis decide, a partir de los datos, si hay evidencia suficiente contra una afirmación sobre la población: que una moneda es equilibrada, que un nuevo modelo no mejora al anterior o que unos recuentos siguen una distribución dada. Este tema presenta el marco de Neyman y Pearson (hipótesis nula y alternativa, región de rechazo, errores de tipo~I y~II, potencia), el $p$-valor, los contrastes $z$ y $t$ para medias y proporciones, el contraste $\chi^{2}$ de bondad de ajuste y el problema de las comparaciones múltiples, con su aplicación a la evaluación de modelos de lenguaje en muchos benchmarks. Se sigue el capítulo~10 de Wasserman~\cite{wasserman}, «Hypothesis Testing and p-values», las lecturas~17b, 18 y~19 del curso 18.05 del MIT~\cite{mit1805}, las clases~5, 7 y~11 de Stats~200~\cite{stats200} y los capítulos~9 y~11 de OpenStax~\cite{openstax}.

\section{Planteamiento de un contraste}

\begin{definicion}[Hipótesis, estadístico y región de rechazo]\label{def:t5-contraste}
Sea $\{F_\theta:\theta\in\Theta\}$ un modelo y $\Theta_0\subset\Theta$. Se contrastan la \emph{hipótesis nula} $H_0:\theta\in\Theta_0$ frente a la \emph{hipótesis alternativa} $H_1:\theta\in\Theta\setminus\Theta_0$. Una hipótesis es \emph{simple} si especifica un único valor de $\theta$ y \emph{compuesta} en otro caso. Un \emph{contraste} consiste en un \emph{estadístico de contraste} $T=T(X_1,\dots,X_n)$ y una \emph{región de rechazo} $R$: si $T\in R$ se rechaza $H_0$; si no, no se rechaza.
\end{definicion}

La hipótesis nula representa la situación por defecto («no hay efecto», «la moneda es equilibrada», «los dos modelos son igual de buenos») y el contraste busca evidencia contra ella. Las alternativas más comunes para un parámetro real son la \emph{bilateral} $H_1:\theta\neq\theta_0$ y las \emph{unilaterales} $H_1:\theta>\theta_0$ o $H_1:\theta<\theta_0$.

\begin{definicion}[Errores, nivel y potencia]\label{def:t5-errores}
Se comete un \emph{error de tipo~I} si se rechaza $H_0$ siendo cierta, y un \emph{error de tipo~II} si no se rechaza siendo falsa:
\begin{center}
\small
\begin{tabular}{@{}lcc@{}}
\toprule
 & $H_0$ cierta & $H_0$ falsa \\
\midrule
No se rechaza $H_0$ & decisión correcta & error de tipo~II \\
Se rechaza $H_0$ & error de tipo~I & decisión correcta \\
\bottomrule
\end{tabular}
\end{center}
El contraste tiene \emph{nivel de significación} $\alpha$ si $\PP_\theta(T\in R)\le\alpha$ para todo $\theta\in\Theta_0$. La \emph{función de potencia} es $\pi(\theta)=\PP_\theta(T\in R)$; para $\theta\in\Theta\setminus\Theta_0$, $\pi(\theta)$ es la \emph{potencia} contra esa alternativa y $\beta(\theta)=1-\pi(\theta)$ es la probabilidad de error de tipo~II~\cite[cap.~10]{wasserman}.
\end{definicion}

Los dos errores no se pueden reducir a la vez con un tamaño muestral fijo: agrandar la región de rechazo reduce los errores de tipo~II y aumenta los de tipo~I. El convenio es fijar $\alpha$ (habitualmente $0{,}05$ o $0{,}01$) y, entre los contrastes de nivel $\alpha$, preferir los más potentes. Las dos hipótesis no son simétricas: no rechazar $H_0$ no prueba que sea cierta, solo que los datos no aportan evidencia suficiente contra ella.

\section{El \texorpdfstring{$p$}{p}-valor}

\begin{definicion}[$p$-valor]\label{def:t5-pvalor}
Supóngase que se rechaza $H_0$ para valores grandes de un estadístico $T$ y que $H_0$ es simple. Si el valor observado es $t_{\mathrm{obs}}$, el \emph{$p$-valor} es
\[
p=\PP_{H_0}(T\ge t_{\mathrm{obs}}),
\]
la probabilidad, calculada suponiendo $H_0$ cierta, de obtener un resultado al menos tan extremo como el observado. (Si se rechaza para valores pequeños, $p=\PP_{H_0}(T\le t_{\mathrm{obs}})$; en un contraste bilateral con estadístico de distribución simétrica, $p=\PP_{H_0}(\abs T\ge\abs{t_{\mathrm{obs}}})$.)
\end{definicion}

\begin{proposicion}[Decisión con el $p$-valor]\label{prop:t5-pvalor}
Sea un contraste que rechaza $H_0$ si $T\ge c_\alpha$, donde $c_\alpha$ cumple $\PP_{H_0}(T\ge c_\alpha)=\alpha$ y $T$ tiene distribución continua bajo $H_0$. Entonces se rechaza $H_0$ al nivel $\alpha$ si y solo si $p\le\alpha$.
\end{proposicion}

\begin{proof}
Sea $G(t)=\PP_{H_0}(T\ge t)$, que es continua y no creciente, de modo que $p=G(t_{\mathrm{obs}})$ y $G(c_\alpha)=\alpha$. Si $t_{\mathrm{obs}}\ge c_\alpha$, entonces $p=G(t_{\mathrm{obs}})\le G(c_\alpha)=\alpha$. Recíprocamente, si $t_{\mathrm{obs}}<c_\alpha$, entonces $p=\PP_{H_0}(T\ge t_{\mathrm{obs}})=\PP_{H_0}(t_{\mathrm{obs}}\le T<c_\alpha)+\alpha$, y el primer sumando es positivo si la densidad de $T$ no se anula entre $t_{\mathrm{obs}}$ y $c_\alpha$, de modo que $p>\alpha$. (Si la densidad se anulara en ese tramo, se toma como $c_\alpha$ el menor valor con $G(c_\alpha)=\alpha$ y el argumento es el mismo.)
\end{proof}

El $p$-valor es, por tanto, el menor nivel al que los datos permiten rechazar $H_0$, y resume en un número la fuerza de la evidencia contra $H_0$: cuanto menor, más incompatibles son los datos con $H_0$.

\begin{observacion}[Lo que el $p$-valor no es]\label{obs:t5-pvalor}
El $p$-valor \emph{no} es la probabilidad de que $H_0$ sea cierta, ni la probabilidad de que los resultados se deban al azar, ni mide el tamaño del efecto. Un $p$-valor muy pequeño puede acompañar a un efecto irrelevante si $n$ es enorme, y un efecto grande puede dar un $p$-valor alto si $n$ es pequeño. Por eso conviene informar siempre del efecto estimado con su intervalo de confianza, además del $p$-valor~\cite[lecturas~17b y~20]{mit1805}.
\end{observacion}

\section{Contrastes \texorpdfstring{$z$}{z} y \texorpdfstring{$t$}{t}}

Los contrastes de esta sección usan como estadístico una estimación tipificada con su error estándar calculado bajo $H_0$. En todos ellos, las regiones de rechazo de nivel $\alpha$ son $T>q_{1-\alpha}$ (alternativa $>$), $T<-q_{1-\alpha}$ (alternativa $<$) y $\abs T>q_{1-\alpha/2}$ (bilateral), donde $q$ es el cuantil de la distribución de referencia.

\begin{center}
\small
\begin{tabular}{@{}llll@{}}
\toprule
Contraste & $H_0$ & Estadístico & Distribución bajo $H_0$ \\
\midrule
$z$ para la media ($\sigma$ conocida) & $\mu=\mu_0$ & $Z=\dfrac{\Xbar-\mu_0}{\sigma/\sqrt n}$ & $\Nor(0,1)$ (datos normales) \\[2ex]
$t$ para la media & $\mu=\mu_0$ & $T=\dfrac{\Xbar-\mu_0}{S/\sqrt n}$ & $t_{n-1}$ (datos normales) \\[2ex]
$z$ para una proporción & $p=p_0$ & $Z=\dfrac{\widehat p-p_0}{\sqrt{p_0(1-p_0)/n}}$ & $\approx\Nor(0,1)$ ($n$ grande) \\
\bottomrule
\end{tabular}
\end{center}

La distribución bajo $H_0$ se sigue del \cref{teo:t3-suma-normales}, del \cref{cor:t3-t} y del \cref{cor:t3-binomial}, respectivamente. En el contraste de una proporción, el error estándar se calcula con $p_0$, no con $\widehat p$, porque se razona bajo $H_0$.

\begin{ejemplo}[Contraste $z$ unilateral]\label{ej:t5-z}
Se afirma que un modelo genera respuestas de $\mu_0=200$ tokens de media, con $\sigma=80$. En $n=64$ respuestas se observa $\xbar=221$. Se contrasta $H_0:\mu=200$ frente a $H_1:\mu>200$ al nivel $0{,}05$. El estadístico vale $z=(221-200)/(80/8)=2{,}1>z_{0{,}95}=1{,}645$, y $p=1-\Phi(2{,}1)=0{,}0179<0{,}05$: se rechaza $H_0$.
\end{ejemplo}

\section{Potencia y tamaño muestral}

\begin{proposicion}[Potencia del contraste $z$ unilateral]\label{prop:t5-potencia}
Para el contraste $z$ de $H_0:\mu=\mu_0$ frente a $H_1:\mu>\mu_0$ al nivel $\alpha$ con datos $\Nor(\mu,\sigma^{2})$, la función de potencia es
\[
\pi(\mu)=1-\Phi\Bigl(z_{1-\alpha}-\frac{(\mu-\mu_0)\sqrt n}{\sigma}\Bigr).
\]
\end{proposicion}

\begin{proof}
Se rechaza si $\Xbar>\mu_0+z_{1-\alpha}\sigma/\sqrt n$. Si la media verdadera es $\mu$, entonces $(\Xbar-\mu)/(\sigma/\sqrt n)\sim\Nor(0,1)$, y
\[
\pi(\mu)=\PP_\mu\Bigl(\frac{\Xbar-\mu}{\sigma/\sqrt n}>\frac{\mu_0-\mu}{\sigma/\sqrt n}+z_{1-\alpha}\Bigr)=1-\Phi\Bigl(z_{1-\alpha}-\frac{(\mu-\mu_0)\sqrt n}{\sigma}\Bigr).\qedhere
\]
\end{proof}

En el \cref{ej:t5-z}, si la media real fuera $220$, la potencia sería $1-\Phi(1{,}645-2)=\Phi(0{,}355)\approx0{,}64$: con $64$ respuestas, el contraste detectaría esa diferencia solo en dos de cada tres ocasiones. El \cref{eje:t5-potencia} deduce el tamaño muestral necesario para alcanzar una potencia dada.

\section{Contraste \texorpdfstring{$\chi^{2}$}{chi cuadrado} de bondad de ajuste}

Sea una variable categórica con $k$ categorías y probabilidades $p_1,\dots,p_k$. En $n$ observaciones independientes se obtienen los recuentos $O_1,\dots,O_k$, con $\sum_jO_j=n$. Se quiere contrastar $H_0:p_j=p_{0j}$ para todo $j$, con $p_{0j}$ dados.

\begin{definicion}[Estadístico $\chi^{2}$ de Pearson]\label{def:t5-chi2}
Sean $E_j=np_{0j}$ los recuentos \emph{esperados} bajo $H_0$. El estadístico es
\[
X^{2}=\sum_{j=1}^{k}\frac{(O_j-E_j)^{2}}{E_j}.
\]
Bajo $H_0$, y si todos los $E_j$ son al menos~$5$, $X^{2}$ sigue aproximadamente una distribución $\chi^{2}_{k-1}$. Se rechaza $H_0$ al nivel $\alpha$ si $X^{2}>\chi^{2}_{k-1;1-\alpha}$~\cite[§11.2]{openstax}.
\end{definicion}

Los $k-1$ grados de libertad se deben a la restricción $\sum_jO_j=n$: conocidos $k-1$ recuentos, el último queda determinado. Si se estiman $m$ parámetros del modelo a partir de los mismos datos, se restan también: $k-1-m$ grados de libertad. Si algún $E_j$ es menor que $5$, se agrupan categorías contiguas.

\begin{ejemplo}[¿Es equilibrado un dado?]\label{ej:t5-dado}
En $120$ lanzamientos de un dado se obtienen los recuentos $25,17,15,23,24,16$ para las caras $1$ a $6$. Bajo $H_0$ (dado equilibrado), $E_j=20$ y
\[
X^{2}=\frac{25+9+25+9+16+16}{20}=\frac{100}{20}=5.
\]
Con $5$ grados de libertad, $\chi^{2}_{5;0{,}95}=11{,}07$; como $5<11{,}07$, no se rechaza $H_0$ al $5\,\%$ (el $p$-valor es aproximadamente $0{,}42$).
\end{ejemplo}

\section{Comparaciones múltiples}

Si se realizan muchos contrastes, cada uno al nivel $\alpha$, la probabilidad de cometer \emph{algún} error de tipo~I crece rápidamente.

\begin{proposicion}\label{prop:t5-multiples}
Si se realizan $m$ contrastes independientes al nivel $\alpha$ y todas las hipótesis nulas son ciertas, la probabilidad de rechazar al menos una es $1-(1-\alpha)^{m}$, y el número esperado de rechazos (todos falsos) es $m\alpha$.
\end{proposicion}

\begin{proof}
Cada contraste no rechaza con probabilidad $1-\alpha$; por independencia, ninguno rechaza con probabilidad $(1-\alpha)^{m}$. El número de rechazos es una suma de $m$ indicadores con esperanza $\alpha$ cada uno.
\end{proof}

Con $m=20$ y $\alpha=0{,}05$, la probabilidad de al menos un «descubrimiento» falso es $1-0{,}95^{20}\approx0{,}64$.

\begin{definicion}[Tasa de error por familia]\label{def:t5-fwer}
La \emph{tasa de error por familia} (FWER, \emph{family-wise error rate}) de un procedimiento que realiza $m$ contrastes es la probabilidad de cometer al menos un error de tipo~I~\cite[clase~11]{stats200}.
\end{definicion}

\begin{teorema}[Corrección de Bonferroni]\label{teo:t5-bonferroni}
Si cada uno de los $m$ contrastes se realiza al nivel $\alpha/m$ (es decir, se rechaza la hipótesis $i$ si $p_i\le\alpha/m$), la tasa de error por familia es a lo sumo $\alpha$, cualquiera que sea la dependencia entre los contrastes y cualquiera que sea el conjunto de hipótesis nulas ciertas.
\end{teorema}

\begin{proof}
Sea $I_0$ el conjunto de índices de las hipótesis nulas ciertas, con $m_0=\abs{I_0}\le m$, y sea $A_i$ el suceso «se rechaza $H_{0,i}$». Para $i\in I_0$, $\PP(A_i)\le\alpha/m$ porque el contraste $i$ tiene nivel $\alpha/m$. Por la desigualdad de la unión (subaditividad de la probabilidad),
\[
\text{FWER}=\PP\Bigl(\bigcup_{i\in I_0}A_i\Bigr)\le\sum_{i\in I_0}\PP(A_i)\le m_0\frac{\alpha}{m}\le\alpha.\qedhere
\]
\end{proof}

La corrección de Bonferroni es sencilla y muy general, pero conservadora: con $m$ grande reduce mucho la potencia. Existen procedimientos más potentes, como el de Holm o el de Benjamini y Hochberg, que controla una cantidad distinta, la tasa de falsos descubrimientos~\cite[clase~11]{stats200}; quedan fuera de este módulo.

\begin{observacion}[Comparar modelos en muchos benchmarks]\label{obs:t5-benchmarks}
La evaluación de modelos de lenguaje es un terreno abonado para las comparaciones múltiples: se comparan dos modelos en decenas de benchmarks, se prueban muchas variantes de \emph{prompt} y se informa de la mejor, o se ordena una tabla con cientos de modelos. Si en $20$ benchmarks independientes no hay ninguna diferencia real, se esperan $20\cdot0{,}05=1$ diferencias «significativas» al $5\,\%$. Para comparar la exactitud de dos modelos se puede usar un contraste $z$ de diferencia de proporciones (\cref{eje:t5-dos-modelos}); cuando ambos modelos responden a las \emph{mismas} preguntas, es preferible un contraste pareado que use las diferencias pregunta a pregunta, porque elimina la variabilidad debida a la dificultad de cada pregunta (ampliación).
\end{observacion}

\begin{observacion}[Errores frecuentes]\label{obs:t5-errores}
(a) Formular la hipótesis nula a partir de los datos ya vistos, o elegir el contraste unilateral después de ver en qué sentido va el efecto.
(b) Concluir «$H_0$ es cierta» porque no se rechaza.
(c) Interpretar el $p$-valor como $\PP(H_0\mid\text{datos})$ (\cref{obs:t5-pvalor}).
(d) En el contraste $\chi^{2}$, usar proporciones o porcentajes en lugar de recuentos, o no restar los parámetros estimados de los grados de libertad.
(e) Realizar muchos contrastes y destacar solo los significativos sin corregir por multiplicidad.
\end{observacion}

\begin{observacion}[Conexión con el curso]
Los módulos~22 y~23 comparan modelos con exactitudes medidas sobre conjuntos finitos de preguntas y discuten la significación de las diferencias, la contaminación de los benchmarks y su saturación; las ideas de este tema (error estándar, potencia, comparaciones múltiples) son las que permiten decidir si una mejora de un punto en una tabla de resultados es real. En el módulo~21, las preferencias humanas entre pares de respuestas se agregan con modelos probabilísticos cuyas diferencias también deben valorarse frente al ruido de muestreo.
\end{observacion}

\begin{lecturas}
\item Wasserman~\cite{wasserman}, cap.~10 «Hypothesis Testing and p-values» (incluye el contraste de Wald, el $\chi^{2}$ de Pearson y los contrastes múltiples).
\item Orloff y Bloom~\cite{mit1805}, lecturas~17b «Null Hypothesis Significance Testing I», 18 «NHST II», 19 «NHST III» y 20 «Comparison of frequentist and Bayesian inference».
\item Stats~200~\cite{stats200}, clases~5 «Testing a simple null hypothesis», 7 «Composite hypotheses and the t-test» y 11 «Testing multiple hypotheses».
\item OpenStax~\cite{openstax}, cap.~9: §9.1 «Null and Alternative Hypotheses», §9.2 «Outcomes and the Type I and Type II Errors» y §9.5 «Additional Information and Full Hypothesis Test Examples»; §11.2 «Goodness-of-Fit Test».
\item Rice~\cite{rice}, cap.~9 «Testing Hypotheses and Assessing Goodness of Fit».
\end{lecturas}

\section*{Ejercicios del tema 5}
\addcontentsline{toc}{section}{Ejercicios del tema 5}

\begin{ejercicio}{1}{eje:t5-vf}
Un contraste da $p=0{,}03$. Indique si cada afirmación es verdadera o falsa. (a) La probabilidad de que $H_0$ sea cierta es $0{,}03$. (b) Al nivel $0{,}05$ se rechaza $H_0$. (c) Al nivel $0{,}01$ se rechaza $H_0$. (d) Si $H_0$ fuera cierta, un resultado al menos tan extremo como el observado ocurriría con probabilidad $0{,}03$. (e) El efecto observado es grande. (f) Si no se hubiera rechazado $H_0$, quedaría demostrado que es cierta.
\end{ejercicio}

\begin{ejercicio}{2}{eje:t5-moneda}
Una moneda se hace girar sobre su canto $250$ veces y sale cara $140$ veces. Sea $\theta$ la probabilidad de cara y $H_0:\theta=0{,}5$ frente a $H_1:\theta\neq0{,}5$. (a) Calcule el $p$-valor con la aproximación normal y corrección por continuidad. (b) ¿Se rechaza $H_0$ al nivel $0{,}10$? ¿Y al $0{,}05$? (c) ¿Cuántas caras harían falta para rechazar al nivel $0{,}01$? (d) Con la región de rechazo de nivel $0{,}05$, estime la potencia contra $\theta=0{,}6$.
\fuente{adaptado de Orloff y Bloom, MIT 18.05~\cite{mit1805}, \emph{problem set}~8, problema~1.}
\end{ejercicio}

\begin{ejercicio}{2}{eje:t5-poligrafo}
En un estudio sobre detectores de mentiras, $140$ personas dicen la verdad y $140$ mienten. De las sinceras, $131$ son clasificadas como sinceras y $9$ como mentirosas; de las que mienten, $15$ son clasificadas como sinceras y $125$ como mentirosas. Con $H_0$: «la persona dice la verdad», describa los errores de tipo~I y~II y estime sus probabilidades, el nivel y la potencia del detector.
\fuente{adaptado de Orloff y Bloom, MIT 18.05~\cite{mit1805}, \emph{problem set}~8, problema~2.}
\end{ejercicio}

\begin{ejercicio}{2}{eje:t5-radar}
Tres radares miden la velocidad de cada coche con errores independientes $\Nor(0,5^{2})$ (en km/h) y se usa la media $\Xbar$ de las tres lecturas. El límite es $40$ y se quiere multar a lo sumo al $4\,\%$ de los conductores que circulan exactamente a $40$. (a) Plantee el contraste y halle el umbral de multa. (b) Si nadie superase el límite, ¿qué fracción de las multas sería errónea? (c) Calcule la potencia contra $\mu=45$. (d) ¿Cuántos radares harían falta para que esa potencia fuera $0{,}9$?
\fuente{adaptado de Orloff y Bloom, MIT 18.05~\cite{mit1805}, \emph{problem set}~8, problema~3.}
\end{ejercicio}

\begin{ejercicio}{2}{eje:t5-uniforme}
Se observa un único valor $x$ de una distribución $\Unif(0,\theta)$ para contrastar $H_0:\theta=2$ frente a $H_1:\theta\neq2$, y se rechaza $H_0$ si $x\le0{,}1$ o $x\ge1{,}9$. (a) Calcule la probabilidad de error de tipo~I. (b) Calcule la probabilidad de error de tipo~II si $\theta=2{,}5$.
\fuente{adaptado de Orloff y Bloom, MIT 18.05~\cite{mit1805}, \emph{problem set}~8, problema~4.}
\end{ejercicio}

\begin{ejercicio}{2}{eje:t5-t-test}
Una muestra normal de tamaño $n=16$ tiene $\xbar=11$ y $s^{2}=4$. Contraste $H_0:\mu=10$ al nivel $0{,}05$ (a) frente a $H_1:\mu\neq10$ y (b) frente a $H_1:\mu>10$. Dé en cada caso el valor crítico y acote el $p$-valor con la \cref{tab:t}.
\fuente{adaptado de Orloff y Bloom, MIT 18.05~\cite{mit1805}, \emph{problem set}~9, problema~1.}
\end{ejercicio}

\begin{ejercicio}{2}{eje:t5-openstax-z}
Un fabricante afirma que sus envases contienen de media $16{,}43$ onzas, con $\sigma=0{,}8$. Una muestra de $15$ envases da $\xbar=16$. Contraste al nivel $0{,}05$ si la media es menor que la anunciada, calculando el estadístico y el $p$-valor.
\fuente{adaptado de OpenStax, \emph{Introductory Statistics}~2e~\cite{openstax}, §9.5, ejemplo~9.14.}
\end{ejercicio}

\begin{ejercicio}{2}{eje:t5-openstax-prop}
Un tratamiento es eficaz, según su fabricante, en el $25\,\%$ de los casos o menos. En una muestra de $42$ casos es eficaz en $17$. Al nivel $0{,}01$, ¿hay evidencia de que la proporción de éxitos es mayor que $0{,}25$?
\fuente{adaptado de OpenStax, \emph{Introductory Statistics}~2e~\cite{openstax}, §9.5, ejemplo~9.19.}
\end{ejercicio}

\begin{ejercicio}{2}{eje:t5-chi-dias}
Las ausencias de $60$ empleados se reparten entre los días laborables así: lunes $15$, martes $12$, miércoles $9$, jueves $9$, viernes $15$. Contraste al nivel $0{,}05$ si las ausencias se reparten uniformemente entre los cinco días.
\fuente{adaptado de OpenStax, \emph{Introductory Statistics}~2e~\cite{openstax}, §11.2, ejemplo~11.2.}
\end{ejercicio}

\begin{ejercicio}{2}{eje:t5-benford}
Según la ley de Benford, el primer dígito $k$ de muchas cantidades contables tiene probabilidades $0{,}301;\allowbreak\ 0{,}176;\allowbreak\ 0{,}125;\allowbreak\ 0{,}097;\allowbreak\ 0{,}079;\allowbreak\ 0{,}067;\allowbreak\ 0{,}058;\allowbreak\ 0{,}051;\allowbreak\ 0{,}046$ para $k=1,\dots,9$. En $100$ partidas de una nota de gastos, los primeros dígitos aparecen $7,13,12,9,9,13,11,10,16$ veces. Contraste al nivel $0{,}001$ si los datos siguen la ley de Benford.
\fuente{adaptado de Orloff y Bloom, MIT 18.05~\cite{mit1805}, \emph{problem set}~9, problema~4.}
\end{ejercicio}

\begin{ejercicio}{2}{eje:t5-dos-modelos}
Dos modelos se evalúan sobre muestras independientes de $1000$ preguntas cada una: A acierta $780$ y B, $750$. (a) Plantee el contraste $H_0:p_A=p_B$ frente a $H_1:p_A\neq p_B$. (b) Bajo $H_0$, estime la varianza de $\widehat p_A-\widehat p_B$ con la proporción combinada y calcule el estadístico $z$ y el $p$-valor. (c) ¿Es significativa la diferencia al $5\,\%$?
\fuente{adaptado de Stats~200, Stanford~\cite{stats200}, \emph{homework}~3, problema~3.}
\end{ejercicio}

\begin{ejercicio}{2}{eje:t5-multiples}
Se compara un modelo con otro en $20$ benchmarks independientes y, en realidad, ambos son igual de buenos en todos. (a) Si cada comparación se contrasta al $5\,\%$, ¿cuál es la probabilidad de encontrar al menos una diferencia significativa? ¿Cuántas se esperan? (b) ¿A qué nivel debe hacerse cada contraste con la corrección de Bonferroni? ¿Cuál es entonces exactamente la probabilidad de al menos un falso positivo?
\end{ejercicio}

\begin{ejercicio}{3}{eje:t5-sidak}
Se realizan $n$ contrastes independientes con $p$-valores $P_1,\dots,P_n$. (a) Si todas las nulas son ciertas y los $P_i$ son \iid\ $\Unif(0,1)$, demuestre que $\PP(\min_iP_i\le t)=1-(1-t)^{n}$ para $t\in(0,1)$. (b) Deduzca que realizar cada contraste al nivel $1-(1-\alpha)^{1/n}$ da una probabilidad de algún falso positivo exactamente igual a $\alpha$. (c) Demuestre que $\alpha/n\le1-(1-\alpha)^{1/n}$, de modo que este procedimiento es al menos tan potente como el de Bonferroni. (\emph{Indicación:} desigualdad de Bernoulli, $(1+x)^{n}\ge1+nx$ para $x\ge-1$.)
\fuente{adaptado de Stats~200, Stanford~\cite{stats200}, \emph{homework}~4, problema~4.}
\end{ejercicio}

\begin{ejercicio}{3}{eje:t5-uniforme-p}
Sea $T$ un estadístico con función de distribución continua y estrictamente creciente $F_0$ bajo $H_0$, y sea $p=1-F_0(T)$ el $p$-valor del contraste que rechaza para valores grandes de $T$. Demuestre que, bajo $H_0$, $p\sim\Unif(0,1)$, y deduzca que $\PP_{H_0}(p\le\alpha)=\alpha$.
\end{ejercicio}

\begin{ejercicio}{3}{eje:t5-dualidad}
Sean datos $\Nor(\mu,\sigma^{2})$ con $\sigma$ conocida. Demuestre que, para cada $\mu_0$, el contraste $z$ bilateral de nivel $\alpha$ \emph{no} rechaza $H_0:\mu=\mu_0$ si y solo si $\mu_0$ pertenece al intervalo de confianza de nivel $1-\alpha$ del \cref{teo:t4-ic-z}. Interprete el resultado.
\fuente{adaptado de Orloff y Bloom, MIT 18.05~\cite{mit1805}, problemas de la clase~23, problema~2.}
\end{ejercicio}

\begin{ejercicio}{3}{eje:t5-potencia}
Con la función de potencia de la \cref{prop:t5-potencia}: (a) demuestre que $\pi$ es creciente en $\mu$ y que $\pi(\mu_0)=\alpha$; (b) demuestre que, para $\mu>\mu_0$ fijo, $\pi(\mu)\to1$ cuando $n\to\infty$; (c) demuestre que la potencia contra $\mu_1=\mu_0+\delta$ ($\delta>0$) es al menos $1-\beta$ si y solo si
\[
n\ge\Bigl(\frac{(z_{1-\alpha}+z_{1-\beta})\,\sigma}{\delta}\Bigr)^{2};
\]
(d) aplíquelo al \cref{ej:t5-z} con $\delta=20$, $\alpha=0{,}05$ y potencia $0{,}9$.
\end{ejercicio}
